%%%%%%%%%%%%%%%%%%%%%%%%%%%%%%%%%%%%%%%%%
% Simple Sectioned Essay Template
% LaTeX Template
%
% This template has been downloaded from:
% http://www.latextemplates.com
%
% Note:
% The \lipsum[#] commands throughout this template generate dummy text
% to fill the template out. These commands should all be removed when 
% writing essay content.
%%%%%%%%%%%%%%%%%%%%%%%%%%%%%%%%%%%%%%%%%

%----------------------------------------------------------------------------------------
%	PACKAGES AND OTHER DOCUMENT CONFIGURATIONS
%----------------------------------------------------------------------------------------

\documentclass[12pt]{book} % Default font size is 12pt, it can be changed here
		\textheight = 28cm
		\textwidth = 19cm
		\topmargin = -1cm
		\oddsidemargin = -1cm
		\parindent = 5mm

\usepackage{geometry} % Required to change the page size to A4
\geometry{a4paper} % Set the page size to be A4 as opposed to the default US Letter

\usepackage{graphicx} % Required for including pictures

\usepackage{float} % Allows putting an [H] in \begin{figure} to specify the exact location of the figure
\usepackage{wrapfig} % Allows in-line images such as the example fish picture

%\usepackage{lipsum} % Used for inserting dummy 'Lorem ipsum' text into the template

\linespread{1.2} % Line spacing

%\setlength\parindent{0pt} % Uncomment to remove all indentation from paragraphs
\usepackage[utf8]{inputenc}
\usepackage[spanish]{babel}
\usepackage[T1]{fontenc}
\usepackage{fancyhdr}

% Configurar los encabezados, pies de pagina y paginas de capitulo
% Encabezados
\lhead[\thepage]{CAPÍTULO \thechapter. \rightmark}
\chead[]{}
\rhead[CAPÍTULO \thechapter. \leftmark]{\thepage}

\renewcommand{\headrulewidth}{0.5pt}

% Pie de pagina

\lfoot[]{\today}
\cfoot[]{}
\rfoot[J.Carlos Ávila]{}
\renewcommand{\footrulewidth}{0pt}

\usepackage{savesym}
\usepackage{amsmath}
\savesymbol{iint}
\usepackage{txfonts}
\restoresymbol{TXF}{iint}


\usepackage[x11names,table]{xcolor}
\usepackage{pstricks}
\usepackage[colorinlistoftodos, textwidth=2cm, shadow]{todonotes}
%\usepackage{hyperref}



\usepackage[colorlinks]{hyperref}
\usepackage[nogroupskip,nopostdot]{glossaries}
\setglossarystyle{altlist}
\makenoidxglossaries




%\usepackage[toc,style=altlistgroup,hyperfirst=false]{glossaries}


\hypersetup{
    colorlinks=true,
    linkcolor=rosa1,
    filecolor=magenta,      
    urlcolor=cyan,
}

\urlstyle{same}

\graphicspath{{./images/}} % Specifies the directory where pictures are stored

\definecolor{miorange}{rgb}{0.11, 0.43, 0.21}
\definecolor{rosa1}{RGB}{236, 46, 80}

%%%%%%%%%%%%%%%%%%%%%%%%%%%%%%%%%%%%%%%%%%%%%%%%%%%%%%%%%%%%%%%%%%%%%%%%%%%%%%%%%%%%%%%%%%%%%%%%%%%%%%%%%%%
%----------------------------------------------------------------------------------------
%	begin {Glosario}
%----------------------------------------------------------------------------------------

\newglossaryentry{SO}{name={SO},description={Es el sistema o conjunto de aplicaciones que permiten que una computadora lleven a cabo sus funciones.}}
\newglossaryentry{Linux}{name={Linux},description={GNU/Linux Sistema Operativo creado y distribuido por miles de 
													personas alrededor del mundo}}
													
\newglossaryentry{Windows}{name={Windows},description={Sistema Operativo creado y distribuido por Microsoft.}}

\newglossaryentry{Mac}{name={Mac},description={Sistema Operativo creado y distribuido por Apple}}

%%%%%%%%%%%%%%%%%%%%%%%%%%%%%%%%%%%%%%%%%%%%%%%%  %%%%%%%%%%%%%%%%%%%%%%%%%%%%%%%%%%%%%%%%%%%%%%%%%%%%%%%%%
\newglossaryentry{}{name={},description={}}
%%%%%%%%%%%%%%%%%%%%%%%%%%%%%%%%%%%%%%%%%%%%%%%%  %%%%%%%%%%%%%%%%%%%%%%%%%%%%%%%%%%%%%%%%%%%%%%%%%%%%%%%%%

%----------------------------------------------------------------------------------------
%	End {Glosario}
%----------------------------------------------------------------------------------------


\pagestyle{fancy}

\begin{document}

\lhead[\thepage]{CAPÍTULO \thechapter. \rightmark}
\rhead[CAPÍTULO \thechapter. \leftmark]{\thepage}


%----------------------------------------------------------------------------------------
%	TITLE PAGE
%----------------------------------------------------------------------------------------

\begin{titlepage}

\newcommand{\HRule}{\rule{\linewidth}{0.5mm}} % Defines a new command for the horizontal lines, change thickness here

\center % Center everything on the page
 
%----------------------------------------------------------------------------------------
%	HEADING SECTIONS
%----------------------------------------------------------------------------------------

\textsc{\LARGE Instituto Tecnológico de San Juan del Río}\\[1.5cm] % Name of your university/college
\textsc{\Large Centro de Investigación en ciencia Aplicada y Tecnología Avanzada}\\[0.5cm] % Major heading such as course name
%\textsc{\large Minor Heading}\\[0.5cm] % Minor heading such as course title

%----------------------------------------------------------------------------------------
%	TITLE SECTION
%----------------------------------------------------------------------------------------

\HRule \\[0.4cm]
{ \huge \bfseries Amplificación interactiva de contenido por medio de la detección de la dirección de la mirada.}\\[0.4cm] % Title of your document
\HRule \\[1.5cm]
 
%----------------------------------------------------------------------------------------
%	AUTHOR SECTION
%----------------------------------------------------------------------------------------

\begin{minipage}{0.4\textwidth}
\begin{flushleft} \large
\emph{Autor:}\\
J. Carlos \textsc{\'Avila Resendiz} % Your name
\end{flushleft}
\end{minipage}
~
\begin{minipage}{0.4\textwidth}
\begin{flushright} \large
\emph{Supervisor:} \\
Dr. Joaquin  \textsc{Salas Rodriguez} % Supervisor's Name
\end{flushright}
\end{minipage}\\[4cm]

% If you don't want a supervisor, uncomment the two lines below and remove the section above
%\Large \emph{Author:}\\
%John \textsc{Smith}\\[3cm] % Your name


%----------------------------------------------------------------------------------------
%	DATE SECTION
%----------------------------------------------------------------------------------------

{\large \today}\\[1cm] % Date, change the \today to a set date if you want to be precise

%----------------------------------------------------------------------------------------
%	LOGO SECTION
%----------------------------------------------------------------------------------------

\includegraphics{./imagenes/itsjr_s.jpg}\\ % Include a department/university logo - this will require the graphicx package

 
%----------------------------------------------------------------------------------------

\vfill % Fill the rest of the page with whitespace
\newpage
$\ $
\thispagestyle{empty}
\end{titlepage}

%----------------------------------------------------------------------------------------
%	TABLE OF CONTENTS
%----------------------------------------------------------------------------------------j

\pagenumbering{roman} 
\tableofcontents % Include a table of contents

\newpage % Begins the essay on a new page instead of on the same page as the table of contents
\thispagestyle{empty} 
%\appendix
%----------------------------------------------------------------------------------------
%	INTRODUCTION
%----------------------------------------------------------------------------------------
\pagenumbering{arabic}	
%\setcounter{page}{1}
	
			 
\newpage		 
\pagestyle{fancy}


\chapter{GENERALIDADES}
\thispagestyle{empty}
\markboth{GENERALIDADES}{GENERALIDADES}

\begin{minipage}{0.5\textwidth}
	\begin{flushleft} \large
	%\emph{•} \\
	\scriptsize	\textsl{\large “El auténtico genio consiste en la capacidad para evaluar información incierta, 
								aleatoria y contradictoria.”}\\
	\scriptsize \textbf{Winston Churchill, estadista.}
	\end{flushleft}
\end{minipage}\\[4cm]			

\newpage
\section{Objetivos}
	\subsection{Objetivo general}
	
		Desarrollar una aplicación de amplificación interactiva para computadoras con sistema operativo Windows, 
		que asista a personas con bajas capacidades visuales, por medio del seguimiento y estimación de la dirección 
		de la mirada sobre la pantalla de la computadora y en base a ello amplificar la zona de la pantalla en la 
		que enfoca la vista.
	
	\subsection{Objetivos específicos}
		\begin{itemize}
			\item Detección precisa y confiable del movimiento del globo ocular, con la ayuda de software de 
			procesamiento de imágenes digitales.
			\item Hacer uso de las API’s del sistema operativo que proveen las herramientas que magnifican la zona de
			la pantalla seleccionada.
			\item Integrar los dos componentes anteriores y de esa forma obtener un magnificador con interacción visual.
			\item Una vez se cuente con un prototipo, realizar pruebas de campo.
		\end{itemize}
		

\newpage
\section{Justificación}
	Pese al avance desmesurado de la tecnología en los últimos años en donde las capacidades de los dispositivos se
	duplica cada cierto tiempo, respondiendo de forma bastante precisa la emblemática ley de 
	\href{https://en.wikipedia.org/wiki/Moore's_law}{Moore} hay aun a día de hoy
	ciertas cuestiones que no han sido abordadas, quizá en gran parte debido al amplio panorama de problemas que se
	pueden afrontar con soluciones tecnológicas y de alguna forma ayudar a solventar o/y hacer más fácil las mismas.
	
	Aun si los programas de asistencia a personas con capacidades diferentes están a día de hoy cobrando mayor 			  	relevancia en prácticamente todos los aspectos sociales, pues en la actualidad las posibilidades de llevar una vida
	productiva y sin las limitaciones de antaño, son ya una realidad, entre las herramientas que se proporcionan a este
	sector de la población están las llamadas tecnologías de asistencia o accesibilidad en entornos informáticos, mismas
	que van desde iconos monocromáticos de un mayor tamaño, hasta lectores de pantalla y lupas, siendo estas últimas el
	principal componente proporcionado por las herramientas de accesibilidad de los Sistemas Operativos \gls{SO} actuales,
	siendo común en los tres mas importantes \gls{Linux}, \gls{Windows}, \gls{Mac}.
	
	Siendo de los tres el segundo, Windows, en el cual se enfocaran los esfuerzos de hacer converger las herramientas de
	accesibilidad ya mencionadas y las tecnologías de visión por computadora CV, para ofrecer a los discapacitados 
	visuales una forma de hacer uso de la tecnología, mismos que según datos de la OMS de 2002, eran mas de 161 millones
	de personas, en especifico de computadoras, sin que su limitante visual les impida el poder interactuar con el 
	equipo.
	
	Específicamente el segmento de la población con discapacidad visual en el que se enfoca el desarrollo de este
	proyecto es el de personas que cuentan con cierto grado de visión, o lo que se conoce como resto visual, pues,
	siempre que exista un resto visual por mínimo que sea se debe potenciar su uso para alcanzar el máximo desarrollo
	posible

\newpage 
\section{Caracterización de la empresa}
	\subsection{Datos generales de la empresa}
	
	\subsection{Descripción del departamento o área de trabajo}

\newpage
\section{Problemas a resolver}

\newpage
\section{Alcances y limitaciones}
	\subsection{Alcances}
	\subsection{Limitaciones}




\chapter{FUNDAMENTACIÓN TEÓRICA}
\markboth{FUNDAMENTACIÓN TEÓRICA}{FUNDAMENTACIÓN TEÓRICA} 
\thispagestyle{empty}

\section{Ingeniería del software}

\section{Herramientas de desarrollo}

\section{Lenguajes de programación}

\section{Metodologías de desarrollo de software}

\section{Eye Gaze }

\chapter{DESCRIPCIÓN DE ACTIVIDADES REALIZADAS}
\markboth{ACTIVIDADES REALIZADAS}{ACTIVIDADES REALIZADAS}
\thispagestyle{empty}
\newpage
\section{Análisis}

\section{Diseño}

\section{Desarrollo}

\section{Pruebas}

\section{Implementación}

\section{Retroalimentación}

\section{Resultados}

\section{Conclusiones y recomendaciones}

\section{Referencias Bibliográficas \& Glosario} 	
%----------------------------------------------------------------------------------------
%	BEGIN GLOSARIO
%----------------------------------------------------------------------------------------			
\newpage


\printnoidxglossaries
\pagenumbering{roman}

%----------------------------------------------------------------------------------------
%	END GLOSARIO
%----------------------------------------------------------------------------------------


%---------------------------------------------------------------------------------------
%	BIBLIOGRAPHY
%----------------------------------------------------------------------------------------

\begin{thebibliography}{99} % Bibliography - this is intentionally simple in this template




%\bibitem{Luckham} 
%David Luckham
%[\textit{The Power of Events - An Introduction to Complex Event Processing in Distributed Enterprise Systems}]. 
%Addison-Wesley, ISBN 0-201-72789-7.


%\bibitem{Tello} 
%Adolfo Lozano Tello
%[\textit{Iniciación a la programación utilizando lenguajes visuales orientados a eventos.}] [ISBN978-84-7897-714-7]
%Ed.Bellisco Ediciones Técnicas y Científicas, ISBN 84-95279-49-5. ISBN 978-84-95279-49-1.


%\bibitem{} 
%Ricardo Barona Vázquez
%[\textit{Asterisk©, telefonía IP en software libre}].
%\href{Asterisk}{ http://www.enterate.unam.mx/artic/2008/abril/art3.html}

\end{thebibliography}


%----------------------------------------------------------------------------------------
 
\end{document}


